\documentclass[12pt,a4paper]{article}

% =========================================================
% Packages
% =========================================================
\usepackage{amsmath,amssymb,mathtools}
\usepackage{geometry}
\geometry{margin=1in}
\usepackage{microtype}

\usepackage{caption}
\usepackage{subcaption}

\usepackage{tikz}
\usetikzlibrary{arrows.meta,calc,angles,quotes,positioning,decorations.markings,patterns,decorations.pathreplacing,shapes.geometric}
\usepackage{tikz-cd}
\usepackage{tikz-3dplot}

\usepackage{enumitem}
\usepackage{float}
\usepackage{xcolor}
\usepackage{hyperref}

\usepackage{pgfplots}
\pgfplotsset{compat=1.18}

\usepackage{fancyhdr}
\usepackage{amsthm}

% =========================================================
% Operators / macros
% =========================================================
\DeclareMathOperator{\ord}{ord}
\DeclareMathOperator{\Tr}{Tr}
\DeclareMathOperator{\PSL}{PSL}
\DeclareMathOperator{\Area}{Area}
\DeclareMathOperator{\Isom}{Isom}
\DeclareMathOperator{\arsinh}{arsinh}
\DeclareMathOperator{\codim}{codim}
\DeclareMathOperator{\Span}{span}

\newcommand{\II}{\mathrm{I\!I}}
\newcommand{\I}{\mathrm{I}} % optional: first fundamental form
\newcommand{\C}{\mathbb{C}}
\newcommand{\R}{\mathbb{R}}
\newcommand{\HH}{\mathbb{H}}
\newcommand{\hyp}{\mathbb{H}}
\newcommand{\DD}{\mathbb{D}}
\newcommand{\T}{\mathsf{T}} % transpose symbol
\newcommand{\sph}{\mathrm{sph}}
\newcommand{\n}{6}

\newcommand{\Solution}{\par\medskip\noindent\textbf{Solution.}\ }
\newcommand{\Theory}{\par\medskip\noindent\textbf{Theory.}\ }
\newcommand{\Proof}{\par\medskip\noindent\textbf{Proof.}\ }

% =========================================================
% Header / footer
% =========================================================
\pagestyle{fancy}
\fancyhf{}
\fancyfoot[C]{\thepage}
\fancyhead[L]{\textit{Differential Geometry MATH3D}}
\fancyhead[R]{\textit{Dev}}

% =========================================================
% Theorems
% =========================================================
\theoremstyle{plain}
\newtheorem{theorem}{Theorem}[section]
\newtheorem{lemma}[theorem]{Lemma}
\newtheorem{proposition}[theorem]{Proposition}
\newtheorem{corollary}[theorem]{Corollary}

\theoremstyle{definition}
\newtheorem{definition}{Definition}
\newtheorem{example}[theorem]{Example}
\newtheorem*{Example}{Example}
\newtheorem{remark}[theorem]{Remark}

% =========================================================
% TikZ style presets (kept from your file)
% =========================================================
\definecolor{PoleRed}{RGB}{214,49,56}
\definecolor{SimpleBlue}{RGB}{31,119,180}
\definecolor{Gold}{RGB}{247,183,51}

\tikzset{
	>={Latex[length=3.4mm]},
	axis/.style={black, line width=1.0pt, ->},
	pole/.style={red!70!black, fill=red!70!black},
	polemark/.style={red!70!black, thick},
	pole2/.style={PoleRed, line width=1.8pt},
	removable/.style={black, line width=1.6pt},
	leader/.style={black!70, line width=0.9pt},
	labelfunc/.style={font=\large},
	labeltext/.style={black!85, font=\normalsize},
	callout/.style={
		draw=black!40, fill=white, rounded corners=3pt,
		inner sep=5pt, outer sep=5pt,
		text=black, opacity=1, text opacity=1
	},
	note/.style={black!70, font=\scriptsize}
}

\tikzset{
	->-/.style={
		postaction={decorate},
		decoration={
			markings,
			mark=at position 0.18 with {\arrow{Stealth[length=2.2mm]}}
		}
	},
	p10loop/.style={thick}
}

% =========================================================
% Poincaré disk helper (kept from your file)
% =========================================================
% Draw a Poincaré disk geodesic between boundary angles a,b (degrees).
% The geodesic is the unique circle arc orthogonal to the unit circle.
% Robust handling of diameter/antipodal case.
\newcommand{\DiskGeodesic}[3][]{%
	\pgfmathsetmacro{\ax}{cos(#2)}%
	\pgfmathsetmacro{\ay}{sin(#2)}%
	\pgfmathsetmacro{\bx}{cos(#3)}%
	\pgfmathsetmacro{\by}{sin(#3)}%
	
	% det = ax*by - ay*bx
	\pgfmathsetmacro{\det}{\ax*\by - \ay*\bx}%
	\pgfmathsetmacro{\detabs}{abs(\det)}%
	
	% diameter / antipodal case: det ~ 0
	\ifdim \detabs pt < 0.000001pt
	\draw[#1] (\ax,\ay) -- (\bx,\by);
	\else
	\pgfmathsetmacro{\h}{(\by-\ay)/\det}%
	\pgfmathsetmacro{\k}{(\ax-\bx)/\det}%
	\pgfmathsetmacro{\Rgeo}{sqrt(\h*\h+\k*\k-1)}%
	\pgfmathsetmacro{\tA}{atan2(\ay-\k,\ax-\h)}%
	\pgfmathsetmacro{\tB}{atan2(\by-\k,\bx-\h)}%
	\pgfmathsetmacro{\dAB}{\tB-\tA}%
	\pgfmathsetmacro{\dAB}{ifthenelse(\dAB>180,\dAB-360,\dAB)}%
	\pgfmathsetmacro{\dAB}{ifthenelse(\dAB<-180,\dAB+360,\dAB)}%
	\pgfmathsetmacro{\tEnd}{\tA+\dAB}%
	\draw[#1] (\h,\k) ++(\tA:\Rgeo)
	arc[start angle=\tA, end angle=\tEnd, radius=\Rgeo];
	\fi
}

% =========================================================
% Title
% =========================================================
\title{Differential Geometry MATH3D}
\author{}
\date{2026-01-24, 16:50}

\begin{document}
	\maketitle
	
	% =========================================================
	% SECTION 1 (Top-level): Geodesic Path on Surface
	% =========================================================
	\section{Geodesic Path on Surface (Mesh vs. Implicit)}
	
	\subsection{Problem: ``edge-walk'' zig-zag artifacts}
	The current path looks like a jagged polyline because it is effectively a \emph{shortest path constrained to mesh edges}
	(i.e., Dijkstra/A* on the vertex-edge graph). On coarse triangle meshes (e.g., Marching Cubes output), this produces
	visible ``staircasing'' and corner kinks.
	
	\subsection{Practical solver options (easy win $\rightarrow$ gold standard)}
	
		\subsubsection*{A* / Dijkstra on a richer graph (cheap upgrade, still graph-based)}
	\textbf{Goal:} reduce zig-zag by allowing paths to cut across faces, not only follow edges.
	
	\paragraph{Add face diagonals / local shortcuts.}
	Add extra edges such as:
	\begin{itemize}
		\item connect vertices across quads (if you can detect them),
		\item connect 2-ring neighbors (neighbors-of-neighbors) with weight equal to straight Euclidean distance.
	\end{itemize}
	\textbf{Pros:} simple, fast. \\
	\textbf{Cons:} still not a true geodesic; can still kink.
	
	\paragraph{Build a ``visibility graph'' on the 1-ring.}
	For each vertex, connect it to all vertices that appear in adjacent faces (not only direct mesh edges).
	\begin{itemize}
		\item This effectively lets A* shortcut within local face neighborhoods.
	\end{itemize}
	\textbf{Pros:} fewer zig-zags. \\
	\textbf{Cons:} can ``cut corners'' incorrectly on coarse meshes.
	
	\medskip
	This is the easiest solver change if minimal math is required.
	
		\subsubsection*{Exact geodesics on triangle meshes (classic discrete geodesics)}
	These compute \emph{exact shortest paths} on polyhedral (triangulated) surfaces.
	
	\paragraph{Chen--Han algorithm.}
	Exact shortest paths on polyhedral surfaces.
	\begin{itemize}
		\item \textbf{Pros:} exact on the mesh surface.
		\item \textbf{Cons:} complex to implement robustly.
	\end{itemize}
	
	\paragraph{Mitchell--Mount--Papadimitriou (MMP).}
	Also exact on polyhedral surfaces.
	\begin{itemize}
		\item \textbf{Pros:} exact.
		\item \textbf{Cons:} even more complex than Chen--Han.
	\end{itemize}
	
		\subsubsection*{Heat Method (best practical solver for apps)}
	This is the \emph{go-to} geodesic distance solver in graphics:
	smooth, robust, and relatively straightforward if sparse linear solves are available.
	
	\paragraph{High-level pipeline.}
	Let $L$ be the Laplace--Beltrami operator (discrete cotan Laplacian) and $M$ the mass matrix.
	\begin{enumerate}
		\item \textbf{Heat diffusion (short time $t$):}
		\[
		(M + tL)\,u = M\,\delta_s
		\]
		where $\delta_s$ is a point source at the start vertex (or barycentric point source).
		
		\item \textbf{Normalize gradient field:}
		\[
		X = -\frac{\nabla u}{\|\nabla u\|}
		\]
		
		\item \textbf{Poisson solve to obtain distance $\phi$:}
		\[
		L\phi = \nabla\cdot X
		\]
		
		\item \textbf{Extract the geodesic curve:}
		backtrace from target to source along $-\nabla\phi$ (integrate on the surface).
	\end{enumerate}
	
	\textbf{Pros:} excellent results, handles noisy/coarse meshes well, widely used. \\
	\textbf{Cons:} requires $L/M$ construction + sparse linear system solves.
	
		\subsubsection*{Fast Marching Method (FMM) on the mesh}
	Compute distance by propagating a front on the triangulated surface.
	\begin{itemize}
		\item \textbf{Pros:} more accurate than edge Dijkstra; classic PDE approach.
		\item \textbf{Cons:} harder than A*; often less pleasant than Heat Method in practice.
	\end{itemize}
	
		\subsubsection*{Intrinsic triangulation / edge flips (high quality, more work)}
	Build an \emph{intrinsic Delaunay triangulation} (IDT) to improve geodesic stability on poor meshes,
	then run geodesic algorithms on the intrinsic mesh.
	\begin{itemize}
		\item \textbf{Pros:} accurate and stable even on low-quality triangulations.
		\item \textbf{Cons:} significant implementation effort.
	\end{itemize}
	
	\subsection{What to do for Marching Cubes implicit surfaces}
	
		\subsubsection*{Immediate fix (short-term)}
	Keep the welded mesh graph, but improve the path quality:
	\begin{itemize}
		\item add face-based start/end nodes (barycentric placement) if not already,
		\item add 2-ring edges to reduce zig-zag,
		\item resample the path to uniform arc-length steps,
		\item Laplacian-smooth the polyline (few iterations),
		\item project smoothed points back onto the implicit surface $f(x,y,z)=0$ (Newton projection),
		\item do this \emph{uniformly across all implicit surfaces} (not only spheres).
	\end{itemize}
	
		\subsubsection*{Real solver (next milestone)}
	Implement the \textbf{Heat Method} on the welded indexed mesh:
	\begin{itemize}
		\item produces a smooth distance field $\phi$,
		\item extracting the path via backtracing yields a clean, stable geodesic,
		\item directly addresses the zig-zag ``edge-walk'' artifact.
	\end{itemize}
	
	\subsection{Geodesics directly on an implicit surface $f(x,y,z)=0$ (ODE formulation)}
	Let
	\[
	S = \{x\in\R^3 : f(x)=0\},\qquad g(x)=\nabla f(x),\qquad n(x)=\frac{g(x)}{\|g(x)\|}.
	\]
	A curve $x(s)\in S$ is a geodesic (parameterized by arc-length $s$) iff:
	\begin{itemize}
		\item it stays on the surface: $f(x(s))=0$,
		\item its acceleration is purely normal: $x''(s) \parallel n(x(s))$,
		\item and $\|x'(s)\|=1$ (arc-length parameterization).
	\end{itemize}
	
		\subsubsection*{Constraint differentiation}
	Differentiate $f(x(s))=0$:
	\[
	\nabla f(x)\cdot x' = 0,
	\]
	and again:
	\[
	x'^{\mathsf T} H_f(x)\,x' + \nabla f(x)\cdot x'' = 0,
	\]
	where $H_f(x)$ is the Hessian matrix of $f$.
	
	Assume geodesic acceleration is normal:
	\[
	x'' = \lambda(s)\,n(x)=\lambda(s)\,\frac{\nabla f(x)}{\|\nabla f(x)\|}.
	\]
	Then
	\[
	\nabla f(x)\cdot x'' = \lambda(s)\,\|\nabla f(x)\|.
	\]
	Therefore
	\[
	\lambda(s) = -\frac{x'^{\mathsf T}H_f(x)\,x'}{\|\nabla f(x)\|}.
	\]
	
		\subsubsection*{Final ODE system (arc-length geodesic)}
	Let $v=x'$ so that $\|v\|=1$. The geodesic ODE becomes:
	\[
	\boxed{
		\begin{aligned}
			x' &= v,\\[4pt]
			v' &= -\frac{v^{\mathsf T}H_f(x)\,v}{\|\nabla f(x)\|^2}\,\nabla f(x),
	\end{aligned}}
	\]
	with constraints
	\[
	f(x)=0,\qquad v\cdot\nabla f(x)=0,\qquad \|v\|=1.
	\]
	
		\subsubsection*{Two-point boundary condition and shooting}
	To connect $p_1$ to $p_2$, choose the initial tangent direction $v(0)$ so that the integrated curve hits $p_2$.
	This is a boundary value problem, typically solved by \textbf{shooting}:
	\begin{enumerate}
		\item Build an orthonormal tangent basis at $p_1$: $(e_1,e_2)\subset T_{p_1}S$ with
		\[
		e_i\cdot n(p_1)=0,\qquad e_1\cdot e_2=0,\qquad \|e_i\|=1.
		\]
		\item Parameterize direction by an angle:
		\[
		v(0;\theta)=\cos\theta\,e_1+\sin\theta\,e_2.
		\]
		\item Integrate to length $L$ to obtain $x(L;\theta)$.
		\item Optimize $(\theta,L)$ to match $p_2$:
		\[
		E(\theta,L)=\|x(L;\theta)-p_2\|^2.
		\]
	\end{enumerate}
	
	\paragraph{Practical issues.}
	\begin{itemize}
		\item Multiple geodesics can exist between two points.
		\item Near the cut locus, the minimization is ill-conditioned.
		\item Numerical drift can push $(x,v)$ off the manifold without stabilization.
	\end{itemize}
	
	\subsection{Stability: keep the curve on the surface}
	ODE integration introduces drift; apply projection steps every step (or every few steps).
	
		\subsubsection*{Newton projection of position onto $f(x)=0$}
	\[
	\boxed{
		x \leftarrow x - \frac{f(x)}{\|\nabla f(x)\|^2}\,\nabla f(x)
	}
	\]
	
		\subsubsection*{Project velocity into the tangent plane and renormalize}
	\[
	\boxed{
		v \leftarrow v - (v\cdot n)\,n,\qquad v\leftarrow \frac{v}{\|v\|}
	}
	\]
	
	\subsection{Analytic gradient $\nabla f$ for implicit surfaces}
	If $f(x,y,z)$ is given as an explicit expression, then
	\[
	\boxed{
		\nabla f(x,y,z)=\left(\frac{\partial f}{\partial x},\frac{\partial f}{\partial y},\frac{\partial f}{\partial z}\right)
	}
	\]
	
		\subsubsection*{Implementation options}
	\begin{enumerate}
		\item \textbf{Symbolic differentiation:} parse $f$ into an AST and build ASTs for $f_x,f_y,f_z$.
		\item \textbf{Automatic differentiation (recommended):} forward-mode AD (dual numbers) with three derivative slots $(dx,dy,dz)$.
		\item \textbf{Finite differences (fallback):}
		\[
		f_x(x,y,z)\approx \frac{f(x+h,y,z)-f(x-h,y,z)}{2h}
		\]
		similarly for $y,z$.
	\end{enumerate}
	
		\subsubsection*{Manual examples}
	If
	\[
	f(x,y,z)=x^2+y^2+z^2-1,
	\quad\text{then}\quad
	\nabla f=(2x,2y,2z),
	\]
	and if
	\[
	f(x,y,z)=x\sin(y)+e^z,
	\quad\text{then}\quad
	\nabla f=\bigl(\sin(y),\,x\cos(y),\,e^z\bigr).
	\]
	
	\subsection{Recommendation}
	\begin{itemize}
		\item \textbf{Shippable solution:} triangulate (Marching Cubes) $\rightarrow$ Heat Method on the mesh $\rightarrow$ backtrace to obtain the path.
		\item \textbf{Experimental ``pure implicit'' mode:} ODE + shooting + projection each step (more fragile but elegant).
	\end{itemize}
	
	% =========================================================
	% SECTION 2 (Top-level): Surface Pipelines / Meshing Strategy
	% =========================================================
	\section{Math3D Surface Meshing Strategy (Implicit / Explicit / Parametric / Weierstrass)}
	
	\subsection{Goal}
	Keep meshing \textbf{robust}, \textbf{predictable}, and \textbf{UI-consistent} across all surface modes in Math3D.
	The key design choice is to \emph{avoid voxelization for smooth parametric surfaces} and reserve CGAL/VTK for implicit
	surfaces where topology is not given by a parameter domain.
	
	\subsection{Parametric-family surfaces (Explicit, Parametric, Weierstrass): do \emph{not} use VTK/CGAL by default}
	For any surface that can be evaluated as a smooth map
	\[
	X(u,v)\in\R^3,
	\]
	\textbf{do not use VTK/CGAL at all} as the default.
	
		\subsubsection*{Parametric mode}
	\[
	X(u,v)\in\R^3.
	\]
	
		\subsubsection*{Explicit mode}
	Explicit graphs are parametric surfaces in disguise:
	\[
	z=f(x,y)
	\quad\Longrightarrow\quad
	X(u,v)=(u,v,f(u,v)).
	\]
	So explicit surfaces reuse the exact same pipeline as parametric surfaces.
	
		\subsubsection*{Weierstrass mode}
	The Weierstrass representation outputs a parametric surface over a domain (with integration / stability checks),
	so the final mesh generation remains parametric:
	\[
	X(u,v).
	\]
	\textbf{Meshing rule: triangulate the parameter domain and map through $X(u,v)$.}
	
	\subsection{Robust default: domain triangulation + adaptive refinement}
	\textbf{Default pipeline for Parametric / Explicit / Weierstrass:}
	\begin{enumerate}[label=\arabic*)]
		\item Start with a coarse triangulation in parameter space $(u,v)$
		(grid triangulation or Delaunay in 2D).
		\item Map vertices via $X(u,v)$ into $\R^3$.
		\item Refine triangles adaptively where needed, based on one or more criteria:
		\begin{itemize}
			\item edge length in $\R^3$ too large,
			\item normal variation across the triangle too large,
			\item curvature estimate too large / rapidly changing,
			\item proximity to known singular regions (optional).
		\end{itemize}
		\item Compute normals / curvature / derived fields as usual.
	\end{enumerate}
	
	\paragraph{Why this is robust.}
	\begin{itemize}
		\item Avoids topology surprises introduced by voxelization.
		\item Gives deterministic control over density exactly where needed (high curvature, seams, singularities).
		\item Works uniformly for Explicit/Parametric/Weierstrass and matches the current renderer model.
	\end{itemize}
	
	\subsection{Implicit surfaces $f(x,y,z)=0$: CGAL is the quality gold standard}
	Implicit surfaces do \emph{not} come with a natural parameter domain:
	\[
	S=\{(x,y,z)\in\R^3 : f(x,y,z)=0\}.
	\]
	Meshing is fundamentally different, and external meshing libraries become the robust choice.
	
		\subsubsection*{Two implicit families}
	
	\paragraph{A) Analytic implicit $f(x,y,z)$ (best quality: CGAL / pygalmesh).}
	If you have an analytic implicit function $f(x,y,z)$, the best mesh quality and robustness usually comes from:
	\begin{itemize}
		\item CGAL implicit surface meshing (via \texttt{pygalmesh} or direct CGAL bindings).
	\end{itemize}
	
	\paragraph{B) Sampled scalar field / voxels / SDF volume (best speed: VTK FlyingEdges / Marching Cubes).}
	If the input is already a discretized volume (grid samples / voxel SDF), the best engineering path is:
	\begin{itemize}
		\item VTK isosurfacing (FlyingEdges3D / Marching Cubes).
	\end{itemize}
	
	\paragraph{Conclusion for implicit meshing.}
	\[
	\textbf{Quality: CGAL} \;>\; \textbf{VTK}
	\qquad\qquad
	\textbf{Speed / preview: VTK} \;>\; \textbf{CGAL}.
	\]
	
	\subsection{One-pipeline rule for predictable UX in Math3D}
	Use a single, easy-to-explain rule:
	
	\subsubsection*{Use param-domain adaptive triangulation for}
	\begin{itemize}
		\item Explicit surfaces ($z=f(x,y)$ as $X(u,v)$),
		\item Parametric surfaces ($X(u,v)$),
		\item Weierstrass surfaces ($X(u,v)$ from the integration pipeline).
	\end{itemize}
	
	\subsubsection*{Use CGAL / pygalmesh for}
	\begin{itemize}
		\item Analytic implicit surfaces $f(x,y,z)=0$ \textbf{when mesh quality matters}.
	\end{itemize}
	
	\subsubsection*{Use VTK isosurfacing for}
	\begin{itemize}
		\item Volumetric / SDF preview meshes,
		\item draft implicit meshes when you already have a 3D grid.
	\end{itemize}
	
	\subsection{Implications for the ``region $\rightarrow$ python $\rightarrow$ return mesh'' workflow}
	Assume the framework can send a selected region and get back a triangle mesh $\{V,F\}$.
	
		\subsubsection*{Parametric / Explicit / Weierstrass}
	\begin{itemize}
		\item Python is not required for baseline meshing.
		\item Mesh can be generated deterministically in JS from the selected parameter region:
		\[
		(u,v)\text{-region} \;\;\text{or}\;\; (x,y)\text{-region}.
		\]
		\item Python may still be useful for heavy adaptive refinement or offline-quality meshing, but it is optional.
	\end{itemize}
	
		\subsubsection*{Implicit}
	\begin{itemize}
		\item Region selection provides:
		\begin{itemize}
			\item bounding box (AABB) in $\R^3$,
			\item target edge length / sizing control,
			\item max triangles / budget.
		\end{itemize}
		\item Python runs CGAL meshing inside the bounding box and returns a mesh:
		\[
		\{V,F\}.
		\]
	\end{itemize}
	
	\subsection{``Mesh quality slider'' mapping (consistent UX)}
	Map a single \emph{Mesh Quality} control to the correct notion of resolution in each mode:
	\begin{itemize}
		\item \textbf{Parametric / Explicit / Weierstrass:} refinement thresholds in parameter-space triangulation
		(edge length, normal variation, curvature variation, refinement depth).
		\item \textbf{Implicit (CGAL quality mode):} sizing field / target edge length / approximation error.
		\item \textbf{Implicit (VTK preview mode):} voxel/grid resolution (and optionally smoothing iterations).
	\end{itemize}
	
	\subsection{Summary}
	\begin{itemize}
		\item \textbf{Do not voxelize} parametric-family surfaces by default: mesh the domain, map through $X(u,v)$, adaptively refine.
		\item \textbf{Use CGAL} for analytic implicit $f(x,y,z)=0$ when you want the best surface quality.
		\item \textbf{Use VTK} for volumetric/SDF preview meshes and fast previews/drafts.
		\item Keep \textbf{one UI mental model}: one quality slider that controls refinement in the appropriate space.
	\end{itemize}
	
	\subsection*{Heat Method Geodesics on a Triangulated Patch (Implementation Notes)}
	
	We assume a triangle mesh patch
	\[
	V=\{p_i\}_{i=1}^N\subset\mathbb{R}^3,\qquad
	F=\{(i,j,k)\}\subset\{1,\dots,N\}^3,
	\]
	with oriented faces \(f=(i,j,k)\). A point \(x\) lying in a face \(f\) is represented by
	\[
	x = b_i p_i + b_j p_j + b_k p_k,
	\qquad
	(b_i,b_j,b_k)\in\Delta^2,\;\; b_i+b_j+b_k=1,\;\; b_\ell\ge 0,
	\]
	i.e.\ by its \emph{barycentric} coordinates in that face.
	
	\subsubsection*{Precompute and cache per mesh}
	
	All quantities below are cached \emph{per mesh} (keyed by a mesh hash) so that repeated geodesic queries reuse the same operators and sparse factorizations.
	
	\paragraph{Mean edge length.}
	Let \(E\) be the set of (undirected) mesh edges. The mean edge length is
	\[
	\bar{h} \;=\; \frac{1}{|E|}\sum_{(a,b)\in E}\|p_a-p_b\|.
	\]
	
	\paragraph{Face areas and normals.}
	For each face \(f=(i,j,k)\), define
	\[
	n_f=(p_j-p_i)\times(p_k-p_i),\qquad
	A_f=\frac{1}{2}\|n_f\|.
	\]
	
	\paragraph{Barycentric basis gradients.}
	Let \(\lambda_i,\lambda_j,\lambda_k\) be the barycentric basis functions on face \(f\).
	Their gradients are constant on \(f\) and can be written as
	\[
	\nabla \lambda_i \;=\; \frac{n_f \times (p_k-p_j)}{2A_f},\qquad
	\nabla \lambda_j \;=\; \frac{n_f \times (p_i-p_k)}{2A_f},\qquad
	\nabla \lambda_k \;=\; \frac{n_f \times (p_j-p_i)}{2A_f}.
	\]
	These three vectors are cached per face and reused to compute gradients of any
	piecewise-linear scalar field.
	
	\paragraph{Cotan Laplacian \(L\).}
	For each undirected edge \((a,b)\) shared by one or two faces, let \(\alpha,\beta\)
	be the angles opposite \((a,b)\) in those incident faces (only \(\alpha\) exists on the boundary).
	Define the cotan weight
	\[
	w_{ab} \;=\; \frac{1}{2}\bigl(\cot\alpha + \cot\beta\bigr)
	\quad\text{(use only the existing angle(s) if the edge is on the boundary).}
	\]
	The cotan Laplacian \(L\in\mathbb{R}^{N\times N}\) is assembled as
	\[
	L_{ab}=-w_{ab}\;\;(a\neq b),\qquad
	L_{aa}=\sum_{b\sim a} w_{ab},
	\]
	where \(b\sim a\) ranges over vertices adjacent to \(a\).
	
	\paragraph{Lumped mass matrix \(M\).}
	Let \(M\) be diagonal with entries obtained by area lumping:
	for each face \(f=(i,j,k)\),
	\[
	M_{ii}\mathrel{+}= \frac{A_f}{3},\qquad
	M_{jj}\mathrel{+}= \frac{A_f}{3},\qquad
	M_{kk}\mathrel{+}= \frac{A_f}{3}.
	\]
	
	\paragraph{Face adjacency.}
	For each face \(f=(i,j,k)\), build the neighbor across each of its three edges.
	Equivalently, store
	\[
	\texttt{nbr}[f][0]=\text{face across edge }(j,k),\quad
	\texttt{nbr}[f][1]=\text{face across edge }(k,i),\quad
	\texttt{nbr}[f][2]=\text{face across edge }(i,j),
	\]
	with value \(-1\) on boundary edges.
	
	\paragraph{Cached sparse factorizations.}
	For fixed \(t\) (defined below) cache a sparse factorization of
	\[
	A \;=\; M + tL.
	\]
	For the Poisson solve cache a factorization of a \emph{pinned} Laplacian \(L_p\),
	where one vertex \(p\) is fixed to remove the nullspace.
	
	\subsubsection*{Heat solve}
	
	\paragraph{Time step.}
	Choose
	\[
	t \;=\; t_{\mathrm{factor}}\,\bar{h}^2.
	\]
	
	\paragraph{Barycentric delta.}
	Let the source point lie in face \(f_s=(i,j,k)\) with barycentric coordinates
	\((b_i,b_j,b_k)\). Form a discrete delta vector \(\delta\in\mathbb{R}^N\) by
	\[
	\delta_i\mathrel{+}=b_i,\qquad
	\delta_j\mathrel{+}=b_j,\qquad
	\delta_k\mathrel{+}=b_k,
	\]
	and \(\delta_\ell=0\) for all other vertices \(\ell\).
	
	\paragraph{Heat equation.}
	Solve for \(u\in\mathbb{R}^N\):
	\[
	(M+tL)\,u \;=\; \delta.
	\]
	
	\subsubsection*{Vector field}
	
	\paragraph{Per-face gradient.}
	For any vertex scalar field \(s\in\mathbb{R}^N\), the gradient on face
	\(f=(i,j,k)\) is
	\[
	\nabla s\big|_f \;=\; s_i\nabla\lambda_i + s_j\nabla\lambda_j + s_k\nabla\lambda_k.
	\]
	Compute \(\nabla u|_f\) for every face.
	
	\paragraph{Normalize.}
	Define the per-face unit vector field
	\[
	X_f \;=\; -\frac{\nabla u|_f}{\|\nabla u|_f\|},
	\]
	with a small \(\varepsilon\) safeguard (if \(\|\nabla u|_f\|<\varepsilon\), set \(X_f=0\)).
	
	\subsubsection*{Poisson solve for distance}
	
	\paragraph{Discrete divergence.}
	Given a face-constant vector field \(X_f\), assemble a vertex RHS \(b\in\mathbb{R}^N\) by
	distributing each face contribution via the barycentric gradients:
	for each face \(f=(i,j,k)\),
	\[
	b_i \mathrel{+}= -A_f\, X_f\cdot\nabla\lambda_i,\qquad
	b_j \mathrel{+}= -A_f\, X_f\cdot\nabla\lambda_j,\qquad
	b_k \mathrel{+}= -A_f\, X_f\cdot\nabla\lambda_k.
	\]
	
	\paragraph{Pinned Laplacian.}
	Since \(L\) has a constant-function nullspace on each connected component,
	pin one vertex index \(p\) (e.g.\ \(p=1\)) by enforcing
	\[
	\phi_p = 0.
	\]
	In practice, replace row \(p\) by the identity constraint:
	\[
	(L_p)_{p\ast}=0,\quad (L_p)_{pp}=1,\quad b_p=0.
	\]
	
	\paragraph{Poisson equation.}
	Solve for \(\phi\in\mathbb{R}^N\):
	\[
	L_p\,\phi \;=\; b.
	\]
	
	\subsubsection*{Geodesic path backtrace (continuous, face-walk)}
	
	Let the target point lie in face \(f_0\) with barycentric coordinates
	\(b^{(0)}=(b^{(0)}_1,b^{(0)}_2,b^{(0)}_3)\). We march a point \(x\) across faces
	following the (negative) gradient of \(\phi\).
	
	\paragraph{Step size.}
	Set
	\[
	h \;=\; \texttt{step\_factor}\cdot \bar{h}.
	\]
	
	\paragraph{Face direction.}
	At the current face \(f=(i,j,k)\), compute
	\[
	g_f \;=\; \nabla\phi\big|_f \;=\; \phi_i\nabla\lambda_i + \phi_j\nabla\lambda_j + \phi_k\nabla\lambda_k,
	\qquad
	d \;=\; -\frac{g_f}{\|g_f\|}.
	\]
	
	\paragraph{Barycentric update.}
	Because barycentric coordinates are linear on each face,
	\[
	\Delta b_\ell \;=\; \nabla\lambda_\ell \cdot (h\,d),\qquad \ell\in\{i,j,k\}.
	\]
	Thus the tentative barycentric coordinates are
	\[
	b'_\ell \;=\; b_\ell + \nabla\lambda_\ell \cdot (h\,d).
	\]
	
	\paragraph{Edge crossing.}
	If all \(b'_\ell\ge 0\), then the new point remains in the same face.
	If some component becomes negative, we intersect the first crossed edge.
	For any \(\ell\) with \(b'_\ell<0\), define the crossing fraction
	\[
	\alpha_\ell \;=\; \frac{b_\ell}{b_\ell - b'_\ell},
	\]
	and take \(\alpha=\min_\ell \alpha_\ell\) (the first boundary hit).
	Update to the edge point
	\[
	b \leftarrow b + \alpha\,(b' - b),
	\]
	so that one barycentric coordinate is \(0\) (on the crossed edge).
	Let \(\ell^\star\) be the index attaining the minimum; then the crossed edge is
	the one \emph{opposite} vertex \(\ell^\star\).
	Move to the adjacent face using the cached adjacency:
	\[
	f \leftarrow \texttt{nbr}[f][\ell^\star].
	\]
	If \(\texttt{nbr}[f][\ell^\star]=-1\), a boundary edge was hit and the march stops (v1).
	
	\paragraph{Monotonicity / tie-break toward the source.}
	To ensure the march moves toward the source level set, accept only steps that decrease \(\phi\)
	(equivalently, move along \(-\nabla\phi\)). If numerical noise yields ambiguity, break ties by
	also choosing the direction that decreases Euclidean distance to the source point.
	
	\paragraph{Stopping.}
	Stop when either
	\[
	\|x - x_s\| < \varepsilon_{\mathrm{pos}}
	\quad\text{or}\quad
	|\phi(x)-\phi(x_s)| < \varepsilon_{\phi}
	\quad\text{or}\quad
	\texttt{steps}>\texttt{max\_steps}.
	\]
	The visited points \(x\) form the geodesic polyline, and its length is the sum of segment lengths.
	
	\paragraph{Caching summary.}
	For each mesh patch we cache:
	\[
	\bigl(L,\; M,\; \{A_f\},\; \{\nabla\lambda\},\; \texttt{nbr},\; \bar{h},\;
	\mathrm{LU}(M+tL),\; \mathrm{LU}(L_p)\bigr),
	\]
	so repeated queries reuse all operators and sparse factorizations.
	
	
\end{document}
